% =====================================================================
%  译校说明（Lead 撰写）
% =====================================================================
\chapter*{译校说明}
\markboth{译校说明}{译校说明}
\addcontentsline{toc}{chapter}{译校说明}

本书是 Edward Frenkel 所著 \emph{Langlands Correspondence for Loop Groups}（Cambridge Studies in
Advanced Mathematics 103, Cambridge University Press, 2007）的中文译注本。以下是体例说明。

\subsection*{一、底本与页码}

\begin{itemize}[leftmargin=1.8em]
    \item 底本为作者 2007 年由 Cambridge University Press 出版的版本，共 406 个 PDF 页。
    \item 原书每章另起新页，PDF 版省略了其中的若干空白页，因此“原书印刷页”与“PDF 页号”的偏移量随章递减：
          PDF p.\,15--139 为 $-14$，p.\,140--176 为 $-13$，p.\,177--231 为 $-12$，p.\,232--260 为 $-11$，
          p.\,261--395 为 $-10$，p.\,396--401 为 $-9$（PDF p.\,1--14 为前置部分，用罗马页码）。例如 PDF 第 15 页为第 1 章第 1 页，PDF 第 315 页为第 10 章第 305 页，PDF 第 389 页为附录第 379 页。
    \item 为避免歧义，源码批注中的页码写作“p.N”，指原书页面图像上印的页码；需要标示扫描件页码时写作“PDF p.NNN”。
    \item 正文、脚注、定理陈述与证明全部译出；参考文献条目照原样以英文转写（作者、标题、期刊、卷期页均保留英文）；
          原书索引基于英文术语与原文页码，本译本以书末附录 B《术语表》代替，不另译索引。
\end{itemize}

\subsection*{二、公式与编号}

\begin{itemize}[leftmargin=1.8em]
    \item 原书的编号公式形如 (2.1-11)，其中 2.1 为节号、11 为该节内的公式序号。
          译本\textbf{逐一沿用原编号}（用 \verb|\tag{2.1-11}| 手工标注），不做重新编号，
          便于与英文原书逐条对读。
    \item 原书中大量显示公式没有编号，译本同样不加编号，保持一一对应。
    \item 正文引用公式时写作“式 (2.1-11)”，不依赖自动交叉引用，以免与原文编号脱节。
    \item 定理、引理、命题、推论、定义、例、注的编号由 \LaTeX{} 环境按节自动生成（形如定理 2.2.2）。
          原书在同一节内偶有重号（如两个命题 3.2.1、两个引理 3.3.1，见 §3.2--§3.3），译本按出现顺序连续编号、
          重复处顺延，并在首次出现处以译注说明原书编号，正文引用统一用译本编号。
    \item 所有公式均已对照原书页面图像逐符号核对；文字层提取仅供定位，不作为公式依据。
\end{itemize}

\subsection*{三、术语}

\begin{itemize}[leftmargin=1.8em]
    \item 术语以中文学界通行译名为准（如“圈群”“顶点代数”“中心扩张”“屏蔽算子”“共形块”）。
    \item 中文学界尚无稳定译名、或保留英文更不易混淆的术语，一律保留英文：
          \emph{oper}、\emph{Miura oper}、\emph{Wakimoto 模}、\emph{Verma 模}、\emph{Weyl 模}、
          \emph{Hecke 本征层}等；首次出现处一般括注英文原词。
    \item 人名与模型名保留英文：Langlands、Kac--Moody、Wakimoto、Verma、Weyl、Hecke、Hitchin、
          Miura、Borel、Cartan、Clifford、Heisenberg、Virasoro、Segal--Sugawara、Feigin--Frenkel、
          Drinfeld--Sokolov、Beilinson--Drinfeld、Satake、Iwahori、Harish--Chandra 等。
    \item 全书术语以书末附录 B《术语表》为唯一裁决书，正文与此表保持一致。
\end{itemize}

\subsection*{四、批注与导读}

\begin{itemize}[leftmargin=1.8em]
    \item 源码中以 \verb|% [Tutor Note: ...]| 标出的 \LaTeX{} 注释是\textbf{译者批注}，编译后不出现在 PDF 中，
          供读者在源码中查阅。批注分四类：
          \begin{enumerate}[leftmargin=1.8em,itemsep=1pt]
              \item \textbf{易错预警}：指标缩并、符号约定、水平与对偶 Coxeter 数、外微分方向等易错处；
              \item \textbf{含糊点拨}：作者以“显然”“容易看出”跳过证明之处，补出所用定理与文献；
              \item \textbf{原文笔误}：如实指出原书印刷错误与正确形式，正文不擅自改动；
              \item \textbf{背景直觉}：该构造在 Langlands 对应与表示论中的位置。
          \end{enumerate}
    \item 每节开头有“本节重点”（keypoints），提炼该节的核心思想、数学工具与关键结论。
    \item 批注不改变正文内容；正文严格照译原文，作者省略的论证不补入正文。
    \item v2 校订版另设\textbf{读者可见的译注}：以“译注：”起头的小号脚注，逐条说明卡点原因、所需背景与
          原文没有的信息，并附来源（原始论文、作者后续论文、nLab/Stacks Project 等）。译注与源码中的
          Tutor Note 并存：前者给读者，后者给研读源码的人。
\end{itemize}

\subsection*{五、v2 校订版说明}

本版为 2026 年的重校订本，相对初版的变化有四项：
\begin{itemize}[leftmargin=1.8em]
    \item 逐章重校正文，订正误译、漏译、术语不一致与中文表达；原书笔误仍照录原文，只用译注指出；
    \item 把原先仅供源码阅读的批注整理、升级为读者可见的译注，并补充外部文献依据；
    \item 在正文前增加《前置知识导引》，在正文后增加《2007 年之后的进展》。两篇均为 \verb|\chapter*|，
          不占原书章节编号，但列入目录；
    \item 统一术语与定理编号，并在书末术语表中登记新术语。
\end{itemize}
两篇新增文字的依据与“已证／猜测／未核实”的区分，见各篇开头的说明。

\subsection*{六、编译}

本译本使用 XeLaTeX{} 编译（文档类 \texttt{ctexbook}，UTF-8 编码）：

\begin{verbatim}
    xelatex main.tex
    xelatex main.tex
    xelatex main.tex
\end{verbatim}

源码结构：\texttt{源码/main.tex} 为主文件，\texttt{源码/preamble.tex} 为导言区，
\texttt{源码/chapters/} 下为各章片段，可直接编辑。

